% Modernised reading — fonds Alexandre Grothendieck, shelfmark 122, pages 1-34.
%
% This is an INTERPRETATION, not a transcription. Everything the critical
% apparatus carried moves into footnotes. In any dispute of substance the
% transcription governs: whoever cites Grothendieck must cite batch-01.fr.tex
% and batch-02.fr.tex.
%
% Pass: Opus 5.5 (claude-opus-5-5), 2026-10-03 — first pass, unchecked against
% the pages by a human. The folder was read whole: two batches, pages 1-20 and
% 21-34, both transcribed. His pages are 2-16 (even), 19, 21-31 (odd), 33, 34;
% the others are blank or typed leaves of other hands, passed by the
% transcription and not described here.
%
% Notation fixed for the whole document.
% — The pages' double-struck R is written R; R Gamma_! is written R Gamma_c.
% — The dimension of the generic manifold of pages 2-10 is d (the pages write
%   n); n is kept for the folder's own n of pages 19-33: W of dimension 2n+1,
%   M (the Milnor fibre) of dimension 2n, W_0 of dimension 2n-1.
% — The gluing identifications of page 12, written alpha, beta on the page,
%   are a, b here: alpha, beta name the components of the restriction
%   R Gamma(P) -> R Gamma(dP) from page 19 on, and that is the use kept.
% — The components of the morphism of triangles of page 21, written
%   (u_0, u_1, u_2) on the page, are (theta_0, theta_1, theta_2): u and u_1
%   name the duality maps of pages 21 and 23.
% — K = R Gamma(M) with its action of pi (the page's R(Gamma, pi)(M), K^.);
%   L = R Gamma(W_0); L^triv with trivial action. The orientation sheaf
%   (a script letter on the page, plain T on p. 2) is or_W. Tate twists are
%   dropped in the topological reading and mentioned once (section on
%   pp. 19-21).
% — Lambda is taken to be a field, or a self-injective ring such as
%   Z/l^nu; the pages identify cohomology of duals with duals of cohomology
%   throughout, which needs it. Supplied, and said in the conventions.
%
% Substantive departures from the pages, each footnoted where it occurs:
% — p. 6: Q = P - P-ring (sic) read as the closure of W - P; the term
%   H^{i+1}(R, F (x) or) of the relative sequence corrected to H^i(R, ...)
%   (Thom isomorphism in codimension 1).
% — p. 14: the homotopy equivalence W - Q -> P holds for any manifold with
%   boundary (collar); it is not what the triviality of dP means. The page's
%   thirteen-vertex diagram is redrawn as the structure of the open book.
% — p. 19 and 23: R Gamma(S^1)-factor written L(-1) on p. 19 and L[-1] on
%   p. 23; both twist and shift belong (L(-1)[-1]). p. 21: Poincaré duality
%   shift [2n] corrected to [-2n], as the page itself has it lower down.
% — p. 21: the class of u lives in H^{2n}(pi, K (x)^L K), not merely in
%   H^{2n}(K (x)^L K), if u is to be a map in D(Lambda[pi]).
% — p. 25: « (c) transposées de (d) » — Verdier duality pairs (d) with (e)
%   and makes (c) self-dual; the overwritten letters are noted, the true
%   pairing stated.
% — p. 27: the boxed sphere criterion needs W a closed connected orientable
%   (2n+1)-manifold; supplied.
% — p. 29: H^i(M) = 0 for i > n does not follow from H^i(W - W_0) = 0 for
%   i > n+1 alone (Wang sequence); it is true (Andreotti-Frankel / affine
%   Lefschetz), which is what the page's bracket « indépendant en fait »
%   asserts.
% — p. 31: with that bound, the « épim si i = n+1 » becomes an isomorphism,
%   and the two outer terms of the boxed sequence vanish; the label « surj »
%   on H^{n-1}(W_0) -> H^n(M)^vee is true only when u vanishes in the middle
%   degree, and is flagged.
% — p. 33: « H^i(V_s) = 0 si i != 0, n-1, n » is false for the link in degree
%   2n-1 (the same page has H^{2n-1}(V_s) = Lambda); read as the statement for
%   the Milnor fibre (p. 31), with the margin's V_xi; flagged as a doubt.
%
% What the folder announces and never establishes: the self-duality of the
% triangles (« vérifier la commutation de R Gamma^pi à R Hom ») is asserted
% and not checked on the pages; the sign of rho' = transpose of rho is left
% with a question mark; the count of « 40, en fait 80 » exact sequences
% breaks off on « mais »; page 34's example breaks off on « Q = ».
\documentclass[11pt,a4paper]{article}
\input{../preamble/grothendieck}

\folder{122}
\pages{1}{34}
\dating{[vers 1968-1971]}
\watermark{Édition de démonstration\\interprétation personnelle de l'œuvre}
\foldertitle{Point singulier isolé de la fibre : notes manuscrites (s.d.) — lecture modernisée du dossier entier}

\begin{document}

\section*{Résumé}

\begin{resume}
Ces feuillets cherchent ce qu'une petite sphère tracée autour d'un point
singulier d'une courbe ou d'une surface permet de dire de ce point.

Prenons l'exemple le plus simple. Dans l'espace des couples de nombres
complexes $(z, w)$, l'équation $z^{2} + w^{3} = 0$ définit une courbe qui a un
point singulier à l'origine : elle y forme une pointe. Entourons l'origine
d'une petite sphère ; dans l'espace des couples complexes, qui a quatre
dimensions réelles, cette sphère en a trois. La courbe la coupe selon un
cercle tordu, qui fait un nœud de trèfle. Toute l'information sur la
singularité est dans la façon dont ce nœud est posé dans la sphère.

La découverte qui organise le sujet, et que ces pages tiennent pour acquise,
est que le complémentaire du nœud est rangé comme un livre ouvert. Le nœud en
est le dos ; les pages sont des surfaces qui s'appuient toutes sur le dos et
qui, en tournant autour de lui, remplissent toute la sphère et reviennent à
leur position de départ. Une page s'appelle aujourd'hui la \emph{fibre de
Milnor}, et le tour complet, qui ramène chaque page sur elle-même mais pas
point par point, la \emph{monodromie}. En dimension supérieure, c'est la même
image : une sphère $W$ de dimension $2n+1$, un « nœud » $W_0$ de dimension
$2n-1$, des pages $M$ de dimension $2n$.

Ce que le dossier fait de cette image est de la cohomologie : une manière de
compter les trous d'un espace, degré par degré. Les feuillets découpent la
sphère en deux morceaux, un épais voisinage du dos et la réunion des pages, et
écrivent toutes les suites exactes qui relient les trous des morceaux à ceux
du tout. Puis ils remarquent que chacune de ces suites est en dualité avec une
autre — à peu près comme, en algèbre linéaire, une application est en dualité
avec sa transposée. Une seule application, notée $u$, porte tout ce qui n'est
pas formel : c'est la forme d'intersection des cycles d'une page, qui dit
comment deux trous de la page se croisent.

La conclusion, quand $W$ est vraiment une sphère, est que presque tout est
déterminé. La sphère n'ayant pas de trous, le nœud et les pages n'en ont qu'en
quelques degrés, et dans les deux degrés du milieu une suite exacte de six
termes, symétrique, relie les trous du nœud, ceux de la page et
l'application $u$. C'est là que vit tout ce que la singularité a de
particulier.

Deux choses donnent à ces pages leur caractère. D'abord, tout y est dit deux
fois : une fois pour des variétés différentiables, où l'on peut dessiner, et
une fois, par un dictionnaire écrit sur une feuille de croquis, pour des
schémas, où le cercle autour duquel tournent les pages devient un point
générique et le tour de monodromie un groupe de Galois. Le dessin sert de
modèle à ce qu'on ne peut pas dessiner. Ensuite, la méthode : avant de choisir,
il recense — quarante suites exactes, « en fait 80 », et la phrase s'arrête
sur « mais ». Le dernier feuillet change de sujet et note un défaut des
conditions de finitude pour des systèmes de modules.

Les noms à chercher ensuite : fibration de Milnor, livre ouvert, suite de
Wang, suite de Gysin, dualité de Poincaré--Lefschetz, théorème de Lefschetz
local, cycles évanescents.

\keywords{Milnor fibration, Milnor fiber, open book decomposition, link of a singularity, isolated hypersurface singularity, monodromy, Poincaré-Lefschetz duality, Mayer-Vietoris sequence, Thom isomorphism, Gysin sequence, Wang sequence, intersection form, local Lefschetz theorem, affine Lefschetz theorem, cohomological purity, vanishing cycles, a-adic systems}
\end{resume}

\pagerange{1}{34}
\section*{Le dossier, sa charpente et ses notations}

Trente-quatre feuillets, dont dix-sept sont de sa main : les pages paires 2 à
16, la page 19, les pages impaires 21 à 31, et les pages 33 et 34. Les autres
sont blancs, ou sont des feuillets dactylographiés d'autres mains dont il a
souvent réutilisé le verso ; la transcription les passe, et cette lecture
aussi.\footnote{L'un d'eux est une annonce imprimée de l'Institut Henri
Poincaré pour des conférences commençant le 25 avril 1968. Elle ne date pas
ses notes, mais elle est compatible avec la datation « [vers 1968-1971] » de
Montpellier, que l'on reprend telle quelle. Les pages ne sont pas numérotées de
sa main ; l'ordre est celui de l'inventaire.}

Le dossier se lit en cinq stations.

\begin{itemize}
\item \emph{Pages 2 à 10.} Préliminaires de topologie : une variété à bord, sa
  dualité, et une variété coupée en deux morceaux le long d'une hypersurface ;
  toutes les suites exactes qui en résultent, et leur appariement par dualité.
\item \emph{Pages 12 à 16.} Le modèle : une variété recollée à partir de deux
  fibrés le long de leur bord commun — un livre ouvert — et le dictionnaire qui
  le relie à la situation algébrique d'un point singulier isolé de la fibre.
\item \emph{Pages 19 à 25.} Les triangles distingués (a) à (e) dans la
  catégorie dérivée, l'application de dualité $u$, et l'énoncé que ces
  triangles sont autotransposés.
\item \emph{Pages 25 à 33.} Le cas où $W$ est une sphère : suite de Gysin,
  Lefschetz local, annulations, et la suite exacte des degrés du milieu ; puis
  la même chose dite en termes de fibres générique et spéciale.
\item \emph{Page 34.} Un autre sujet : les conditions de finitude pour les
  systèmes inductifs $\mathfrak{a}$-adiques.
\end{itemize}

\textbf{Les objets.} $W$ est une variété compacte de dimension $2n+1$ — dans
le bon cas une sphère —, $W_0 \subset W$ une sous-variété compacte de
codimension $2$, $Q$ un voisinage tubulaire fermé de $W_0$, isomorphe à
$W_0 \times D^{2}$, et $P$ l'adhérence de $W - Q$, de sorte que
\[
  W = P \cup Q, \qquad P \cap Q = \partial P = \partial Q \simeq W_0 \times S^{1},
\]
et que l'inclusion $P \subset W - W_0$ est une équivalence d'homotopie. Le
morceau $P$ est fibré sur le cercle $S^{1}$, de fibre une variété à bord $M$
de dimension $2n$, de bord $\partial M = W_0$. On note $\pi$ le groupe
fondamental du cercle, qui agit sur la cohomologie de $M$ par monodromie, et
\[
  K = R\Gamma(M, \Lambda) \in D(\Lambda[\pi]), \qquad
  L = R\Gamma(W_0, \Lambda) \in D(\Lambda),
\]
$L^{\mathrm{triv}}$ désignant $L$ muni de l'action triviale de $\pi$. Dans la
situation algébrique que le modèle traduit, $W$ est le bord d'une petite boule
autour du point singulier, $W_0$ son intersection avec la fibre spéciale (le
\emph{link}), $M$ la fibre de Milnor.\footnote{Le dossier écrit
$\mathbb{R}(\Gamma, \pi)(M)$ et $K^{\bullet}$ pour notre $K$, $L^{\bullet}$
pour $L$, $\mathbb{R}\Gamma_{!}$ pour notre $R\Gamma_{c}$, et marque
l'intérieur d'une variété à bord d'un petit rond ($\mathring{P}$,
$\mathring{M}$), ce que l'on garde.}

\textbf{Deux homonymies à ne pas confondre.} La lettre $n$ désigne aux pages 2
à 10 la dimension d'une variété quelconque, et à partir de la page 19 l'entier
tel que $\dim W = 2n+1$ ; on écrit $d$ pour la première. Les lettres $\alpha$,
$\beta$ désignent page 12 deux identifications de recollement, et à partir de
la page 19 les deux composantes de la restriction
$R\Gamma(P) \to R\Gamma(\partial P)$ ; on garde ce second usage, qui est celui
de tout le reste du dossier, et l'on écrit $a$, $b$ pour le premier.

\textbf{Les coefficients.} $\Lambda$ est un anneau de coefficients, que l'on
suppose être un corps, ou plus généralement un anneau auto-injectif comme
$\mathbf{Z}/\ell^{\nu}$, et $(-)^{\vee} = R\mathrm{Hom}_{\Lambda}(-, \Lambda)$
désigne le dual.\footnote{Hypothèse suppléée. Le dossier passe constamment de
la dualité des complexes à celle des groupes de cohomologie,
$H^{i}(K^{\vee}) = H^{-i}(K)^{\vee}$, ce qui demande que
$\mathrm{Hom}(-, \Lambda)$ soit exact. Sur $\mathbf{Z}$ il faudrait ajouter des
termes $\mathrm{Ext}^{1}$. Les cohomologies en jeu sont de type fini, les
variétés étant compactes ; c'est ce que dit la marge de la page 21 : « $K$
parfait relativement à $\Lambda$ ».} On suppose les variétés orientées quand
il le faut ; c'est le cas de toutes celles qui viennent de la géométrie
complexe.

\pagerange{2}{4}
\section*{Variétés à bord et dualité (pages 2 à 4)}

Soit $P$ une variété topologique à bord, purement de dimension $d$, de bord
$\partial P$ et d'intérieur $\mathring{P} = P - \partial P$, et $F$ un système
local de $\Lambda$-modules de type fini. L'inclusion $\mathring{P} \to P$ est
une équivalence d'homotopie, donc
\[
  H^{*}(P, F) \xrightarrow{\ \sim\ } H^{*}(\mathring{P}, F),
  \qquad \text{c'est-à-dire} \qquad H^{*}_{\partial P}(P, F) = 0,
\]
et, en localisant, les faisceaux de cohomologie locale
$\underline{H}^{*}_{\partial P}(F)$ sont nuls : au voisinage d'un point du
bord, $P$ est un demi-espace. La suite exacte de la paire $(P, \partial P)$
s'écrit
\[
  \cdots \to H^{i}_{c}(\mathring{P}, F) \to H^{i}(P, F) \to H^{i}(\partial P, F)
  \to H^{i+1}_{c}(\mathring{P}, F) \to \cdots
\]
et la dualité de Poincaré sur la variété sans bord $\mathring{P}$ donne
\[
  H^{i}_{c}(\mathring{P}, F) \simeq
  H^{d-i}(\mathring{P}, F^{\vee} \otimes \mathrm{or})^{\vee} \simeq
  H^{d-i}(P, F^{\vee} \otimes \mathrm{or})^{\vee},
\]
où $\mathrm{or}$ est le faisceau d'orientation.\footnote{La page note $T$ le
faisceau d'orientation et $\check{F}$ le système local dual. En marge : « si on
travaille avec faisceaux loc. constants […], i.e. fibres de dim finie ». La
page dit « paracompacte », le préfixe étant ajouté au-dessus de la ligne et
entouré. Pour la suite on a besoin que $P$ soit \emph{compacte} : c'est ce qui
rend $H^{i}_{c}(\mathring{P}) \simeq H^{i}(P, \partial P)$ et la cohomologie
de type fini ; hypothèse suppléée.} On obtient ainsi un homomorphisme
\[
  \varphi : H^{d-i}(P, F^{\vee} \otimes \mathrm{or})^{\vee} \longrightarrow H^{i}(P, F),
  \qquad \text{c'est-à-dire} \qquad
  \varphi \in H^{i}(P, F) \otimes H^{d-i}(P, F^{\vee} \otimes \mathrm{or}),
\]
et si $\partial P = \emptyset$ c'est un isomorphisme : c'est la dualité de
Poincaré pour $P$.

Le reste de ces deux pages remplace le bord par un collier. Si
$\widetilde{\partial P} \simeq \partial P \times [0, 1]$ est un épaississement
du bord et $P' = \overline{P - \widetilde{\partial P}}$, alors $P - P'$ se
rétracte sur $\partial P$, et l'excision $(P', \partial P') \to (P, P - P')$
donne
\[
  H^{i}_{P'}(P, F) \simeq H^{i}(P, \partial P; F) \simeq H^{i}_{c}(\mathring{P}, F)
  \simeq H^{i}_{c}(\mathring{P}', F).
\]
Autrement dit, la cohomologie à support compact de l'intérieur se lit aussi
comme une cohomologie à support dans un compact $P'$ que l'on a éloigné du
bord ; c'est la forme sous laquelle elle reviendra.\footnote{Tout le haut de
la page 4 est barré ; on n'en retient que la conclusion, qui est reprise
proprement au-dessous. Les « mod » des formules
$H^{i}(P \bmod \partial P, F)$, qui désignent la cohomologie relative, sont
des lectures douteuses.}

\pagerange{6}{10}
\section*{Une variété coupée en deux (pages 6 à 10)}

Soit $W$ une variété compacte purement de dimension $d$, et $P \subset W$ une
sous-variété à bord de même dimension, adhérence de son intérieur. Soit $Q$
l'adhérence de $W - P$.\footnote{La page écrit $Q = P - \mathring{P}$, ce qui
serait le bord de $P$ ; c'est un lapsus, et la suite de la page n'a de sens
que pour l'adhérence du complémentaire.} Alors $Q$ est aussi une sous-variété à
bord, et
\[
  W = P \cup Q, \qquad R := P \cap Q = \partial P = \partial Q, \qquad
  W - R = \mathring{P} \amalg \mathring{Q},
\]
avec $R$ purement de dimension $d-1$.

\subsection*{Les trois suites et leurs duales}

\emph{La suite de $(W, R)$.} La suite de cohomologie à support dans le fermé
$R$, jointe à l'isomorphisme de Thom en codimension $1$,
$H^{i+1}_{R}(W, F) \simeq H^{i}(R, F \otimes \mathrm{or}_{R/W})$, s'écrit
\[
  \cdots \to H^{i}(W, F) \to H^{i}(P, F) \oplus H^{i}(Q, F)
  \to H^{i}(R, F \otimes \mathrm{or}_{R/W}) \to H^{i+1}(W, F) \to \cdots
\]
puisque $H^{i}(W - R) = H^{i}(\mathring{P}) \oplus H^{i}(\mathring{Q}) =
H^{i}(P) \oplus H^{i}(Q)$. L'orientation normale $\mathrm{or}_{R/W}$ est ici
triviale, $R$ bordant $P$.\footnote{La page écrit le troisième terme
$H^{i+1}(R, F \otimes \mathcal{T}_{R})$ ; l'isomorphisme de Thom en
codimension $1$ abaisse le degré d'une unité, et le terme correct est
$H^{i}(R, \ldots)$, ce qu'exige aussi la comparaison avec Mayer--Vietoris qui
suit. La bulle en marge dit bien « faisceau d'orientation normal […] ici
$\simeq \mathbf{Z}$ ».} Il conjecture — « n'est sans doute autre que » — que
c'est la suite de Mayer--Vietoris du recouvrement $W = P \cup Q$, et c'est le
cas, à la convention de signe près sur la flèche
$H^{i}(P) \oplus H^{i}(Q) \to H^{i}(R)$. Sa forme à supports compacts,
\[
  \cdots \to H^{i}_{c}(\mathring{P}, F) \oplus H^{i}_{c}(\mathring{Q}, F)
  \to H^{i}_{c}(W, F) \to H^{i}_{c}(R, F) \to \cdots,
\]
en est la transposée par dualité de Poincaré, puisque
$H^{i}_{c}(\mathring{P}, F) \simeq H^{d-i}(P, F^{\vee} \otimes
\mathrm{or}_{W})^{\vee}$ et de même pour $Q$.\footnote{La page écrit
$H^{i}_{!}(W - T, F)$ pour $H^{i}_{c}(W - R, F)$ : la lettre $T$ y tient lieu
de $R$, comme dans le tableau récapitulatif de la page 10.}

\emph{La suite de $(W, P)$.} Si $\widetilde{P}$ est un petit épaississement de
$P$ dans $W$, l'excision et la page 4 donnent
$H^{i}_{P}(W, F) \simeq H^{i}_{P}(\widetilde{P}, F) \simeq
H^{i}_{c}(\mathring{P}, F)$, et $H^{i}(W - P) = H^{i}(\mathring{Q}) \simeq
H^{i}(Q)$. La suite de cohomologie à support dans $P$ devient donc
\[
  \cdots \to H^{i}_{c}(\mathring{P}, F) \to H^{i}(W, F) \to H^{i}(Q, F) \to \cdots,
  \qquad H^{i}_{c}(\mathring{P}, F) \simeq H^{d-i}(P, F^{\vee} \otimes \mathrm{or}_{W})^{\vee},
\]
et symétriquement en échangeant $P$ et $Q$.

\subsection*{Le tableau de la page 10}

Les trois suites s'entrelacent autour de $H^{i}(W, F)$ en un diagramme en
tresse : la suite de $(W, P)$, celle de $(W, Q)$ et celle de Mayer--Vietoris
partagent deux à deux leurs termes.

\begin{tikzcd}[column sep=tiny, row sep=small, nodes={font=\scriptsize}]
 & H^{i-1}(R) \arrow[dl] \arrow[dr] & \\
H^{i}_{c}(\mathring{P}) \simeq H^{d-i}(P)^{\vee} \arrow[r] & H^{i}(W) \arrow[dl] \arrow[dr] & H^{d-i}(Q)^{\vee} \simeq H^{i}_{c}(\mathring{Q}) \arrow[l] \\
H^{i}(P) \arrow[r] & H^{i}(R) & H^{i}(Q) \arrow[l]
\end{tikzcd}

Les coefficients sont omis ; les termes duaux sont pris à
coefficients $F^{\vee} \otimes \mathrm{or}_{W}$.\footnote{La page dessine
l'hexagone sans pointe sur le trait de $H^{i-1}(R)$ vers le terme de gauche ;
on l'oriente comme le cobord de la suite de $(Q, R)$, ce qu'il est.
L'hexagone de la page 16 répète celui-ci, prolongé d'un cran vers
$H^{i+1}(W)$.}

Il en tire un tableau, qu'il appelle « Récapitulation ». Notons $S(W, P)$ la
suite de $(W, P)$ ci-dessus, $S_{c}(W, P)$ sa forme à supports compacts et
$MV$ la suite de Mayer--Vietoris. On a
\[
  \left.\begin{aligned}
    \Phi &= S(W, P) \simeq S_{c}(W, Q) \\
    \Phi' &= S(W, Q) \simeq S_{c}(W, P)
  \end{aligned}\right\}\ \text{duales,}
  \qquad
  \left.\begin{aligned}
    \Psi &= S(W, R) \simeq MV(W; P, Q) \\
    \Psi' &= S_{c}(W, R) \simeq MV_{c}(W; P, Q)
  \end{aligned}\right\}\ \text{duales.}
\]
Les identifications $S(W, P) \simeq S_{c}(W, Q)$ tiennent parce que $W$, $P$,
$Q$ sont compacts. $\Psi$ fait intervenir $H(P)$, $H(Q)$, $H(R)$, $H(W)$ ;
$\Phi$ fait intervenir $H(P)^{\vee}$, $H(Q)$, $H(W)$ ; $\Phi'$ fait intervenir
$H(P)$, $H(Q)^{\vee}$, $H(W)$. Il les dispose en losange,
$\Psi$ en haut, $\Psi'$ en bas, $\Phi$ et $\Phi'$ de part et
d'autre.\footnote{Un bloc écrit dans l'autre sens de la feuille, sur la suite
de $(W, Q)$, est barré ; il porte un « ? » sur l'isomorphisme
$H^{i}_{Q}(W) \simeq H^{i}_{c}(W - P)$, qui est juste (c'est la page 8, avec
$P$ et $Q$ échangés). Le tableau encadré qui précède la récapitulation dit la
même chose sous une autre disposition ; les $R$ de sa première ligne sont des
lectures douteuses.}

\pagerange{12}{16}
\section*{Le modèle : un livre ouvert (pages 12 à 16)}

\subsection*{La construction}

Soient $S$, $T$ deux variétés différentiables compactes, $P \to S$ un fibré de
fibre une variété à bord $M$, et $Q \to T$ un fibré de fibre une variété à
bord $N$. Alors $\partial P$ est fibré sur $S$ de fibre $\partial M$, et
$\partial Q$ sur $T$ de fibre $\partial N$. On se donne des difféomorphismes
\[
  a : S \simeq \partial N, \qquad b : T \simeq \partial M
\]
et des trivialisations $\sigma : \partial P \simeq S \times \partial M$,
$\tau : \partial Q \simeq T \times \partial N$.\footnote{La page note
$\alpha$, $\beta$ ce que l'on note $a$, $b$ ; voir les conventions.} On en
déduit
\[
  \partial P \simeq \partial N \times \partial M, \qquad
  \partial Q \simeq \partial M \times \partial N,
\]
d'où, en échangeant les facteurs, un difféomorphisme
$\varphi : \partial P \simeq \partial Q$, et $W$ est la variété obtenue en
recollant $P$ et $Q$ le long de leurs bords via $\varphi$. Les inclusions
\[
  W - Q = P - \partial P \hookrightarrow P, \qquad
  W - P = Q - \partial Q \hookrightarrow Q
\]
sont des équivalences d'homotopie.\footnote{La phrase de la page 14 est en
partie illisible, et semble dire que c'est la trivialité de $\partial P$ qui
« signifie » ces équivalences d'homotopie. Elles tiennent en fait pour toute
variété à bord, par le théorème du collier ; ce que la trivialité des bords
fournit est le difféomorphisme $\varphi$, donc le recollement lui-même.}

Le cas qui sert est celui où $N = D^{2}$ est un disque, de bord $S = S^{1}$, et
où $T = \partial M$. Alors $Q = T \times D^{2}$ est un voisinage tubulaire de
$T$ (l'âme $T \times \{0\}$ étant notre $W_0$), $P$ est fibré sur le cercle de
fibre $M$, et leur bord commun est
\[
  S^{1} \times T = \partial N \times \partial M = \partial P = \partial Q = P \cap Q .
\]
C'est la structure qu'on appelle aujourd'hui un \emph{livre ouvert} : $T$ est
le dos (la reliure), les fibres de $P \to S^{1}$ sont les pages, et chacune a
pour bord le dos.\footnote{Le nom est de Winkelnkemper (1973), postérieur à ces
pages ; la chose, pour le bord d'une petite boule autour d'une singularité
isolée d'hypersurface complexe, est le théorème de fibration de Milnor
(\emph{Singular points of complex hypersurfaces}, 1968). Le dossier ne cite
personne. Il pose le modèle sans démonstration : c'est un cadre, pas un
théorème.}

\begin{tikzcd}[column sep=small, row sep=small]
M \arrow[r, hook] & P \arrow[d] & S^{1} \times T \arrow[l, hook'] \arrow[r, hook] \arrow[dl] \arrow[dr] & Q \arrow[d] & N \arrow[l, hook'] \\
 & S = S^{1} & & T = \partial M &
\end{tikzcd}

Les flèches horizontales sont les inclusions d'une fibre et du bord
commun ; les flèches obliques sont les deux projections du produit
$S^{1} \times T$.\footnote{Le diagramme de la page 14 a treize sommets : outre
ceux d'ici, $W$ (avec les inclusions de $P$ et $Q$), les points-bases $s$, $t$
et $(s, t)$ des bases, et des copies de $\partial N = S^{1}$ et de
$\partial M = T$ placées à côté de $N$ et de $M$. Plusieurs de ses têtes de
flèche sont douteuses ; on n'en reproduit que la structure, qui ne dépend pas
d'elles. La page écrit « $N = D_{1} \sim e$ », lecture douteuse : $N$ est un
disque, contractile ($e$ est le point), et c'est un disque de dimension $2$
puisque son bord est un cercle.}

\subsection*{Le dictionnaire (page 16)}

La page 16 est une feuille de croquis. Elle porte la clef de tout le dossier :
la correspondance entre les morceaux du modèle et les objets de la géométrie
algébrique.
\[
  \begin{array}{c|cccc|cccc}
    \text{modèle} & W & P & Q & R & M & N & T = \partial M & S = \partial N \\ \hline
    \text{géométrie} & V' & V_{\eta} & V'_{s} & V'_{s} \times \eta
      & X'_{\bar\eta} & e & V'_{s} & \eta \sim S^{1}
  \end{array}
\]
Ici $V'$ est un voisinage épointé du point singulier, $V'_{s}$ sa trace sur la
fibre spéciale et $V_{\eta}$ sa trace sur la fibre générique ; le point
générique $\eta$ du disque épointé joue le rôle du cercle, et la fibre
géométrique générique celui de la page du livre. Le groupe $\pi$ de la page
19 est, dans cette lecture, le groupe d'inertie, qui agit sur la fibre de
Milnor par monodromie.\footnote{La page écrit $X'_{\eta}$ sous $P$ et à nouveau
sous $M$. Le modèle exige que ce soient deux objets différents : $P$ est fibré
sur le cercle, $M$ est sa fibre. On lit le second comme la fibre
\emph{géométrique} générique $X'_{\bar\eta}$, ce que la page ne marque pas.
Elle écrit aussi $T = \dot{M}$, $S = \dot{N}$, le point marquant le bord ; le
$N$ de $\dot{N}$ peut se lire $M$.}

La feuille porte encore deux diagrammes en perspective, en forme de cube, l'un
sur les objets du modèle, l'autre sur leurs correspondants, qu'on ne reproduit
pas ; un second exemplaire de l'hexagone de la page 10, prolongé jusqu'à
$H^{i+1}(W)$ ; l'accouplement
$H^{i}(Q) \times H^{d-i-1}(P) \to H^{d-1}(R) \simeq \Lambda$ ; et un
dénombrement. Les dix parties à trois éléments de $\{1, \ldots, 5\}$ sont
écrites en colonnes, des sommes partielles sont encadrées — $20$, $14$, $2$,
$4$ — et le total est entouré : « 40 suites exactes de coh. relative et autant
pour l'homologie… En fait 80 suites exactes, mais ». La phrase s'arrête
là.\footnote{La page écrit
$\binom{5}{2} = \frac{5.4.3}{12.} = 10$ ; le résultat $10$ est juste
($\binom{5}{2} = \binom{5}{3}$, et c'est bien dix triplets qui sont listés),
le calcul intermédiaire ne l'est pas tel qu'écrit. On ne reconstitue pas ce qui
est dénombré au-delà de ce que les inscriptions disent.}

\pagerange{19}{21}
\section*{Les triangles de la fibre de Milnor (pages 19 à 21)}

On revient aux notations des conventions : $\dim W = 2n+1$, $W_0 = T$ le dos,
$M$ la page, $K = R\Gamma(M)$ dans $D(\Lambda[\pi])$, $L = R\Gamma(W_0)$. Toutes
les identifications ci-dessous sont celles du modèle : $R\Gamma(P) \simeq
R\Gamma(W - W_0)$, $R\Gamma(Q) \simeq R\Gamma(W_0)$, et
$R\Gamma_{c}(\mathring{P}) \simeq R\Gamma_{c}(W - W_0)$.\footnote{La dernière
tient parce que $(W - W_0) - \mathring{P} = Q - W_0$ est isomorphe à
$\partial Q \times {]0, 1]}$, dont la cohomologie à support compact est
nulle.} La page 19 écrit quatre triangles distingués, la page 21 un
cinquième.\footnote{Sur la page les triangles sont dessinés, deux de leurs
côtés en pointillé, et le sommet qui reçoit la flèche de degré $1$ est mis
entre crochets et entouré au crayon. On les écrit ici en ligne.}

\textbf{(a)}, dans $D(\Lambda[\pi])$ : la suite de la paire $(M, \partial M)$,
avec son action de monodromie,
\[
  R\Gamma_{c}(\mathring{M}) \longrightarrow K \xrightarrow{\ \rho\ } L^{\mathrm{triv}}
  \xrightarrow{\ +1\ } .
\]
L'action sur $L$ est triviale parce que la monodromie peut être choisie égale à
l'identité près du bord de $M$ — c'est la trivialisation $\sigma$ de la page
12. La flèche $\rho$ est la restriction au bord, « compatible avec cup ».

\textbf{(b)} : on applique le foncteur des invariants dérivés
$R\Gamma^{\pi} = R\mathrm{Hom}_{\Lambda[\pi]}(\Lambda, -)$. Comme $P$ est fibré
sur le cercle de fibre $M$, $R\Gamma^{\pi}(K) \simeq R\Gamma(P)$ ; et
\[
  R\Gamma^{\pi}(L^{\mathrm{triv}}) \simeq R\Gamma(W_0 \times S^{1})
  \simeq L \otimes^{\mathbf{L}}_{\Lambda} R\Gamma(S^{1}) \simeq L \oplus L[-1] .
\]
On obtient la suite de la paire $(P, \partial P)$, avec
$\partial P = S^{1} \times W_0$ :
\[
  R\Gamma_{c}(\mathring{P}) \longrightarrow R\Gamma(P)
  \xrightarrow{\ (\alpha, \beta)\ } L \oplus L[-1] \xrightarrow{\ +1\ } .
\]
La composante $\alpha : R\Gamma(P) \to L$ est la restriction à une copie
$W_0 \times \{\mathrm{pt}\}$ du bord, compatible aux cup-produits. La
composante $\beta : R\Gamma(W - W_0) \to L[-1]$ s'identifie en cohomologie,
dit la marge, au cobord de la suite de cohomologie relative qui relie
$H^{i}(W - W_0)$ et $H^{i-1}(W_0)$ : c'est le \emph{résidu} le long de
$W_0$.\footnote{La page 19 écrit le second facteur $R\Gamma(W_0)(-1)$, la page
23 $L^{\bullet}[-1]$. Les deux appartiennent à la formule : dans le cadre
étale que le modèle traduit, $R\Gamma(\pi, \Lambda) \simeq \Lambda \oplus
\Lambda(-1)[-1]$ pour la partie modérée de l'inertie, et le facteur est
$L(-1)[-1]$ ; dans le modèle topologique la torsion de Tate est invisible. On
l'omet désormais, comme dans l'isomorphisme de Thom
$R\Gamma_{W_0}(W) \simeq L(-1)[-2]$ qui suit. La décomposition
$R\Gamma(S^{1}) \simeq \Lambda \oplus \Lambda[-1]$ est légitime : un complexe
dont la cohomologie est $\Lambda$ en degrés $0$ et $1$ est scindé, l'obstruction
vivant dans $\mathrm{Ext}^{2}_{\Lambda}(\Lambda, \Lambda) = 0$.}

\textbf{(c)} : on applique à (a) le foncteur d'oubli
$D(\Lambda[\pi]) \to D(\Lambda)$, et l'on retrouve la suite de $(M, \partial M)$
sans action :
\[
  R\Gamma_{c}(\mathring{M}) \longrightarrow R\Gamma(M)
  \longrightarrow R\Gamma(W_0) \xrightarrow{\ +1\ } .
\]
\footnote{L'exposant que porte ici $K$ pour marquer l'oubli de l'action est lu
« $\omega$ » (ou « w ») à la page 19 et « ens » au triangle (c') de la page 23 ;
les deux lectures sont douteuses et l'on ne tranche pas. On écrit $K$ pour
l'objet de $D(\Lambda[\pi])$ et $R\Gamma(M)$ pour son image dans
$D(\Lambda)$.}

\textbf{(d)} : la suite de cohomologie relative de $(W, W - W_0)$ — « au choix
pour $(W, P)$ » —, jointe à l'isomorphisme de Thom
$R\Gamma_{W_0}(W) \simeq R\Gamma(W_0)[-2]$ :
\[
  R\Gamma(W_0)[-2] \longrightarrow R\Gamma(W) \longrightarrow R\Gamma(W - W_0)
  \xrightarrow{\ \beta\ } R\Gamma(W_0)[-1] .
\]
La flèche de degré $1$ de ce triangle est, note la marge, « le $\beta$ de
(b) ».\footnote{La page écrit le sommet $R\Gamma_{!}(W)$ ; $W$ étant compacte,
c'est $R\Gamma(W)$. Elle fait aboutir $\beta$ au sommet $R\Gamma(W_0)[-2]$, le
décalage d'une unité étant celui de la flèche de degré $1$ ; on l'écrit
explicitement.} C'est la suite de Gysin.

\textbf{(e)} (page 21) : la suite de la paire $(W, Q)$, par excision
$H(W, Q) \simeq H_{c}(\mathring{P})$ :
\[
  R\Gamma_{c}(\mathring{P}) \longrightarrow R\Gamma(W)
  \xrightarrow{\ \rho_1\ } R\Gamma(Q) \simeq R\Gamma(W_0) \xrightarrow{\ +1\ } .
\]
\footnote{La page étiquette $\beta'$ la flèche de degré $1$ de (e) ; la même
lettre désigne à la page 23 une autre flèche, celle de (b'). Le long de
$\rho_1$ : « compatible avec […] », illisible.}

Les restrictions $W \to P$ et $Q \to \partial P$ définissent un morphisme de
triangles $(\theta_0, \theta_1, \theta_2)$ de (e) dans (b) : $\theta_0$ est
l'identité de $R\Gamma_{c}(\mathring{P})$ (c'est l'isomorphisme $u_0$ que la
page dit induit), $\theta_1 = \rho_0 : R\Gamma(W) \to R\Gamma(P)$ est la
restriction, et $\theta_2 : L \to L \oplus L[-1]$ est l'inclusion du premier
facteur, image inverse par la projection $W_0 \times S^{1} \to W_0$. D'où la
relation
\[
  \rho_1 = \alpha \circ \rho_0
\]
que la page note en marge.\footnote{La page écrit le morphisme
$(u_0, u_1, u_2)$, avec « $u_1 = \rho_0$, $u_2 = \mathrm{id}_1$ » ; on lit
$\mathrm{id}_1$ comme l'identité sur le premier facteur. On renomme ces
composantes, $u$ et $u_1$ désignant plus loin d'autres flèches. La
justification de « $u_0$ isomorphisme » est illisible ; c'est l'excision. Une
note en oblique dans la marge gauche, qui mentionne (a'), est presque
entièrement illisible.}

\pagerange{21}{25}
\section*{Dualité : l'application $u$ et les triangles autotransposés (pages 21 à 25)}

\subsection*{L'application $u$}

Supposons $W$ non singulière, c'est-à-dire variété topologique ; comme $W$ est
déjà une variété au voisinage de $W_0$, cela revient à demander que
$\mathring{P}$, ou $\mathring{M}$, le soit.\footnote{Dans la situation
algébrique, $W$ est toujours lisse près de $W_0$ : la fibre spéciale n'ayant
qu'un point singulier isolé, le morphisme est lisse le long de ses autres
points. La non-singularité de $W$ est donc une hypothèse sur l'espace total au
point singulier.} La dualité de Poincaré sur la variété orientée $\mathring{M}$
de dimension $2n$ donne, de façon équivariante,
\[
  R\Gamma_{c}(\mathring{M}) \simeq K^{\vee}[-2n]
  \qquad \text{dans } D(\Lambda[\pi]),
\]
où $\pi$ agit sur $K^{\vee}$ par la contragrédiente.\footnote{La page écrit
d'abord $[2n]$, crochet repassé, puis $[-2n]$ plus bas ; c'est $[-2n]$ qui est
juste : $H^{i}(K^{\vee}[-2n]) = H^{2n-i}(M)^{\vee} \simeq
H^{i}_{c}(\mathring{M})$. Dans le cadre étale s'ajoute une torsion $(n)$.} La
première flèche de (a) devient alors
\[
  \boxed{\ u : K^{\vee}[-2n] \longrightarrow K \quad \text{dans } D(\Lambda[\pi])\ }
\]
et (a) devient le triangle
\[
  \text{(a')} \qquad K^{\vee}[-2n] \xrightarrow{\ u\ } K \xrightarrow{\ \rho\ }
  L^{\mathrm{triv}} \xrightarrow{\ \rho'\ } K^{\vee}[-2n+1] .
\]
En cohomologie, $u$ est l'application $H^{i}_{c}(\mathring{M}) \to H^{i}(M)$
qui oublie le support ; en degré $n$ c'est la \emph{forme d'intersection} de
la fibre de Milnor.

Comme $K$ est parfait sur $\Lambda$, une flèche $K^{\vee}[-2n] \to K$ est la
même chose qu'une classe de degré $2n$ de $K \otimes^{\mathbf{L}}_{\Lambda} K$ ;
pour qu'elle soit une flèche de $D(\Lambda[\pi])$, la classe doit être prise
dans la cohomologie de $\pi$ :
\[
  u \in H^{2n}\bigl(\pi,\ K \otimes^{\mathbf{L}}_{\Lambda} K\bigr) =
  \mathrm{Hom}_{D(\Lambda[\pi])}(K^{\vee}[-2n], K).
\]
\footnote{La page écrit $\xi \in H^{2n}(K \otimes^{\mathbf{L}}_{\Lambda} K)$,
sans $\pi$, et la lettre $\xi$ est douteuse. Telle quelle, c'est la classe de
l'application sous-jacente dans $D(\Lambda)$ ; l'application de $D(\Lambda[\pi])$
en est un relèvement, et c'est lui que le reste du dossier utilise.}

Il affirme que $\rho$ et $\rho'$ sont transposés l'une de l'autre, avec un
« signe ? », c'est-à-dire que le triangle (a') est
\emph{autotransposé} : son dual, décalé de $-2n$, est le même triangle tourné.
C'est la dualité de Poincaré--Lefschetz pour la variété à bord $M$ : elle
échange $R\Gamma(M)$ et $R\Gamma_{c}(\mathring{M})$, et, $\partial M = W_0$
étant une variété fermée de dimension $2n-1$, elle identifie $L^{\vee}[-2n]$ à
$L[-1]$. Le transposé de $\rho : K \to L$ est alors une flèche
$L[-1] \to K^{\vee}[-2n]$, c'est-à-dire $\rho'$.\footnote{Le signe que la page
interroge dépend des conventions de rotation des triangles et de
l'identification $L^{\vee} \simeq L[2n-1]$ ; on ne le fixe pas. Le même
phénomène donne, sur le degré moyen, ${}^{t}u = (-1)^{n} u$ : la forme
d'intersection d'une variété de dimension $2n$ est symétrique pour $n$ pair,
alternée pour $n$ impair.}

\subsection*{Les triangles (c') et (b')}

Le foncteur d'oubli donne
\[
  \text{(c')} \qquad R\Gamma(M)^{\vee}[-2n] \longrightarrow R\Gamma(M)
  \longrightarrow R\Gamma(W_0) \xrightarrow{\ +1\ },
\]
« encore un triangle autotransposé ». Le foncteur $R\Gamma^{\pi}$ donne
\[
  \text{(b')} \qquad R\Gamma(P)^{\vee}[-2n-1] \xrightarrow{\ u_1\ }
  R\Gamma(P) \xrightarrow{\ (\alpha, \beta)\ } L \oplus L[-1]
  \xrightarrow{\ (\alpha', \beta')\ } ,
\]
autotransposé lui aussi, avec $R\Gamma(P) \simeq R\Gamma(W - W_0)$ et
$R\Gamma(P)^{\vee}[-2n-1] \simeq R\Gamma_{c}(\mathring{P})$ : c'est la
dualité de Poincaré--Lefschetz pour la variété à bord $P$, de dimension
$2n+1$.

Le décalage $-2n-1$ et non $-2n$ est ce que la marge encadrée demande de
vérifier : « la commutation de $R\Gamma^{\pi}$ à $R\mathrm{Hom}_{\Lambda}(-,
\Lambda)$ ». Elle a lieu avec un décalage : pour $\pi = \mathbf{Z}$, groupe de
dualité de Poincaré de dimension $1$ (c'est la dualité sur le cercle),
\[
  R\Gamma^{\pi}(K^{\vee}) \simeq R\Gamma^{\pi}(K)^{\vee}[-1] .
\]
\footnote{Vérification suppléée : la page la demande et ne la fait pas. Dans le
cadre étale, la même formule vaut pour la partie modérée de l'inertie, avec
une torsion $(-1)$. Après (b') la page écrit « [exprimer comment ?] » : c'est
cette commutation qui l'exprime.}

\subsection*{Les données nouvelles}

Il en tire deux structures nouvelles. La première est
\[
  \boxed{\ u_1 : R\Gamma(W - W_0)^{\vee}[-2n-1] \longrightarrow R\Gamma(W - W_0)\ },
\]
l'image de $u$ par $R\Gamma^{\pi}$. La seconde est une paire d'accouplements,
obtenus en transposant $\alpha$ et $\beta$ et en identifiant
$L^{\vee} \simeq L[2n-1]$ par la dualité de Poincaré sur $W_0$ :
\[
  \boxed{\begin{aligned}
    \lambda &: R\Gamma(W_0) \otimes R\Gamma(W - W_0) \longrightarrow \Lambda[-2n], \\
    \mu &: R\Gamma(W_0) \otimes R\Gamma(W - W_0) \longrightarrow \Lambda[-(2n-1)],
  \end{aligned}}
\]
soit, en cohomologie,
\[
  \lambda : H^{2n-i}(W_0) \times H^{i}(W - W_0) \to \Lambda, \qquad
  \mu : H^{2n-1-i}(W_0) \times H^{i}(W - W_0) \to \Lambda .
\]
Par la dualité de Poincaré sur $W_0$, $H^{2n-i}(W_0)^{\vee} \simeq
H^{i-1}(W_0)$ et $H^{2n-1-i}(W_0)^{\vee} \simeq H^{i}(W_0)$, de sorte que ces
accouplements sont des applications
\[
  \lambda^{i} : H^{i}(W - W_0) \to H^{i-1}(W_0), \qquad
  \mu^{i} : H^{i}(W - W_0) \to H^{i}(W_0),
\]
et il le dit : « je dis que c'est $\beta$ », « je dis que c'est $\alpha$ ».
C'est bien ce qu'ils sont, par construction.\footnote{La page écrit les
accouplements avec un $\times$ ; ils sont définis sur le produit tensoriel
dérivé. Les exposants des deux lignes en cohomologie sont récrits par-dessus
d'autres, illisibles ; ceux que l'on donne sont ceux que la page porte en
dernier, et ce sont ceux qu'impose le degré des accouplements.}

La page 25 résume : on a deux données remarquables, $u$ et l'accouplement
\[
  \boxed{\ R\Gamma(W_0)^{\mathrm{triv}} \otimes K \longrightarrow \Lambda[-(2n-1)]\ },
\]
qui n'est autre que $\rho$ : par dualité de $W_0$, une flèche
$K \to L^{\vee}[-(2n-1)] \simeq L^{\mathrm{triv}}$ est la même chose qu'un tel
accouplement. Et « enfin, les triangles (a) et (b) sont transposés », chacun
étant son propre transposé.\footnote{La phrase dit « transposés », sans
préciser ; (a) et (b) vivent dans des catégories différentes et en dimensions
différentes ($2n$ et $2n+1$), et ne peuvent être transposés l'un de l'autre. On
la lit, comme les triangles (c') et (b') de la page 23, au sens de
l'autotransposition.}

\subsection*{Sans hypothèse de non-singularité}

Une remarque (« NB ») étend tout cela : avec des coefficients généraux,
localement constants au voisinage de $W_0$, et les coefficients « duaux » de
l'autre côté, on n'a plus besoin que $W$ soit non singulière. La dualité de
Verdier remplace celle de Poincaré, et elle apparie les triangles comme suit :
(a), (b) et (c) sont chacun transposés d'un triangle du même type, et (d) est
transposé de (e). En effet, le dual de (d) décalé de $-(2n+1)$ est, terme à
terme,
\[
  R\Gamma_{c}(W - W_0) \longrightarrow R\Gamma(W) \longrightarrow R\Gamma(W_0)
  \xrightarrow{\ +1\ },
\]
qui est (e).\footnote{La page écrit : « les suites exactes du type (a), (b),
(e) sont transposées de suites exactes du même type, et les suites exactes du
type (c) sont transposées de celles du type (d) ». La troisième lettre de la
première liste est surchargée et « (e) » est une lecture douteuse. Telle que
transcrite, l'affirmation ne tient pas : (c) est la suite d'une variété à bord
de dimension $2n$, son dual est une suite du même type, tandis que le dual de
(d) est (e), comme on vient de le voir. L'appariement que l'on donne est celui
qui est vrai ; il est compatible avec la lecture « (a), (b), (c) … (e) … (d) »
que les surcharges laissent ouverte.}

\pagerange{25}{31}
\section*{Quand $W$ est une sphère (pages 25 à 31)}

\subsection*{La suite de Gysin}

Supposons que $W$ soit une sphère cohomologique de dimension $2n+1$ :
$H^{i}(W) = \Lambda$ pour $i = 0, 2n+1$, et $0$ sinon. La suite de Gysin (d),
\[
  \cdots \to H^{i-2}(W_0) \to H^{i}(W) \to H^{i}(W - W_0)
  \xrightarrow{\ \beta\ } H^{i-1}(W_0) \to H^{i+1}(W) \to \cdots,
\]
donne aussitôt :

\begin{itemize}
\item[a)] si $i \neq 0, 2n$, le résidu $\beta : H^{i}(W - W_0) \to
  H^{i-1}(W_0)$ est un isomorphisme ;\footnote{Il écrit d'abord la condition
  « $i, i+1 \neq 0, 2n+1$ », puis : « en fait, comme on sait (b), il suffit que
  $i \neq 0, 2n$ ». Pour $i = 2n+1$ les deux membres sont nuls, $P$ étant une
  variété compacte à bord non vide de dimension $2n+1$ et $W_0$ de dimension
  $2n-1$.}
\item[b)] dans les degrés du haut, la portion
  \[
    0 = H^{2n}(W) \to H^{2n}(W - W_0) \hookrightarrow H^{2n-1}(W_0)
    \twoheadrightarrow H^{2n+1}(W) = \Lambda \to H^{2n+1}(W - W_0) = 0 ,
  \]
  avec $H^{2n}(W_0) = 0$. Si $W_0$ est connexe, $H^{2n-1}(W_0) \simeq \Lambda$
  par dualité, et donc $H^{2n}(W - W_0) = 0$ et $H^{2n-1}(W_0) \simeq
  H^{2n+1}(W)$.\footnote{La bulle de la page dit « si $W_0$ est connexe, [ce
  qui est] le cas si $n \geq 2$ », et une autre renvoie à « Lefschetz
  [affine ?] local plus bas » : la connexité de $W_0$ pour $n \geq 2$ résulte
  en effet de l'isomorphisme $H^{0}(W) \simeq H^{0}(W_0)$ du théorème de
  Lefschetz local ci-dessous.}
\end{itemize}

La réciproque est encadrée en marge, et elle est juste : \emph{si $W$
est une variété fermée connexe orientable de dimension $2n+1$, c'est une sphère
cohomologique si et seulement si $\beta : H^{i}(W - W_0) \to H^{i-1}(W_0)$ est
un isomorphisme pour $i \neq 0, 2n$ et injectif pour $i = 2n$} ; quand $W_0$
est connexe, la dernière condition équivaut à $H^{2n}(W - W_0) =
0$.\footnote{Hypothèse suppléée : l'encadré ne dit pas que $W$ est une variété
fermée connexe. Sans elle, les conditions sur $\beta$ ne donnent pas
$H^{0}(W) = H^{2n+1}(W) = \Lambda$ ; avec elle, la suite de Gysin donne
l'annulation de $H^{i}(W)$ pour $1 \leq i \leq 2n$, et la dualité fait le
reste. L'encadré ajoute : « si $n \geq 2$ et $W$ non singulière, il faut que
$H^{2n}(W - W_0) = 0$, $W_0$ connexe ». La page 33 reprend le critère sous
une forme complète, avec les conditions en degrés $0$ et $2n$.}

\subsection*{Lefschetz local}

Supposons l'espace total $X$ régulier au point singulier $x$ de la fibre. Alors
le \emph{théorème de Lefschetz local} donne
\[
  H^{i}(W) \longrightarrow H^{i}(W_0) \quad
  \text{isomorphisme pour } i \leq n-2, \ \text{injectif pour } i = n-1,
\]
et, transposé par les dualités de $W$ (dimension $2n+1$) et de $W_0$
(dimension $2n-1$), il devient l'énoncé sur l'application de Gysin
\[
  H^{i}(W_0) \longrightarrow H^{i+2}(W) \quad
  \text{isomorphisme pour } i \geq n+1, \ \text{surjectif pour } i = n .
\]
\footnote{Les bornes des deux énoncés sont récrites par-dessus d'autres
chiffres, et la seconde, $n$, est douteuse ; celles que l'on donne sont celles
que la dualité impose ($j \leq n-2$ équivaut à $i = 2n-1-j \geq n+1$). La page
écrit $H^{i}_{!}(W_0)$ ; $W_0$ est compacte. Topologiquement, c'est le
théorème de Milnor : le link d'une singularité isolée d'hypersurface de
dimension $n$ est $(n-2)$-connexe. Dans le cadre étale, c'est la forme
cohomologique des énoncés de Lefschetz locaux de SGA~2, et le fait que $W$
soit une sphère cohomologique quand $X$ est régulier est la pureté
cohomologique — connue pour $X$ lisse sur un corps (SGA~4, XVI), démontrée en
général par Gabber bien plus tard. La page ne cite aucun de ces textes.}

Donc la cohomologie de $W$ détermine celle de $W_0$, sauf dans les deux degrés
$n-1$ et $n$ — « complémentaires, i.e. où les cohomologies sont duales pour
$W_0$ ». Et l'on a
\[
  H^{i}(W - W_0) = H^{i}(P) = 0 \qquad \text{pour } i > n+1 .
\]

\subsection*{Les annulations}

Revenant au cas où $W$ est une sphère, on trouve
\[
  H^{i}(W_0) = 0 \quad \text{si } i \neq 0, n-1, n, 2n-1, \qquad
  H^{i}(W - W_0) = 0 \quad \text{si } i \neq 0, 1, n, n+1,
\]
la première par Lefschetz local et sa transposée, la seconde par le résidu
$\beta$.\footnote{La page écrit la seconde annulation « si $i > n+1$ » puis,
entre crochets, « en fait si $i \neq 0, 1, n, n+1$ » ; c'est la forme
crochetée que l'on retient, qui contient l'autre. Le degré $2n$, où $\beta$
n'est pas un isomorphisme, est couvert par b) ci-dessus, pour $n \geq 2$.} La
suite de Gysin est alors triviale, sauf les morceaux
\[
  \begin{aligned}
    &0 \to H^{0}(W) \to H^{0}(W - W_0) \to 0, \qquad
     0 \to H^{1}(W - W_0) \to H^{0}(W_0) \to 0, \\
    &0 \to H^{n}(W - W_0) \to H^{n-1}(W_0) \to 0, \qquad
     0 \to H^{n+1}(W - W_0) \to H^{n}(W_0) \to 0 \quad (n \neq 1), \\
    &0 \to H^{2n-1}(W_0) \to H^{2n+1}(W) \to 0,
  \end{aligned}
\]
et, pour $n = 1$, où les degrés $n$ et $2n-1$ se confondent, le dernier
devient
$0 \to H^{2}(W - W_0) \to H^{1}(W_0) \to H^{3}(W) \simeq \Lambda \to 0$.
Pour $n = 1$, $W_0$ est un entrelacs de la sphère de dimension $3$, qui peut
avoir plusieurs composantes : c'est l'exemple du nœud de trèfle.

Il en déduit, encadré,
\[
  \boxed{\ H^{i}(M, \Lambda) = 0 \quad \text{pour } i > n\ }
\]
et ajoute, entre crochets, que c'est « indépendant en fait » de la
non-singularité de $W$ et de $W_0$. L'énoncé est vrai, et c'est la remarque
entre crochets qui le justifie : $M$ a le type d'homotopie d'un complexe
cellulaire de dimension $n$.\footnote{La borne encadrée est surchargée ; $n$
est douteux, $n+1$ possible. C'est $n$ qui est vrai. Il ne découle pas de la
seule annulation de $H^{i}(W - W_0)$ pour $i > n+1$ : par la suite de Wang
$H^{i}(P) \to H^{i}(M) \xrightarrow{\gamma - 1} H^{i}(M) \to H^{i+1}(P)$, où
$\gamma$ est la monodromie, celle-ci donne seulement que $\gamma - 1$ est
bijectif sur $H^{i}(M)$ pour $i > n$, ce qui n'interdit pas $H^{i}(M) \neq 0$.
L'annulation elle-même est un théorème de type Lefschetz affine : topologiquement,
Andreotti--Frankel et Milnor ($M$ est une variété de Stein de dimension
complexe $n$) ; dans le cadre étale, le théorème de Lefschetz affine d'Artin
(SGA~4, XIV) et ses variantes. C'est d'ailleurs le mot que porte, avec un
point d'interrogation, la marge de la page 27. On remarquera que la page 31
est écrite comme si l'on ne savait que l'annulation pour $i > n+1$.}

\subsection*{La suite exacte des degrés du milieu}

Lorsque $W$ est non singulière, le triangle (c') donne la suite exacte longue
\[
  \cdots \to H^{2n-i}(M)^{\vee} \xrightarrow{\ u\ } H^{i}(M) \to H^{i}(W_0)
  \to H^{2n-i-1}(M)^{\vee} \xrightarrow{\ u\ } H^{i+1}(M) \to H^{i+1}(W_0) \to \cdots
\]
Avec l'annulation précédente, elle donne des isomorphismes en dehors du milieu :
\[
  \begin{aligned}
    H^{i}(M) &\xrightarrow{\ \sim\ } H^{i}(W_0) && \text{pour } i \leq n-2,
      && \text{injectif pour } i = n-1, \\
    H^{i}(W_0) &\xrightarrow{\ \sim\ } H^{2n-i-1}(M)^{\vee} && \text{pour } i \geq n+1 .
  \end{aligned}
\]
\footnote{La page donne le second isomorphisme pour $i \geq n+2$ seulement, avec
« épim si $i = n+1$ », et note « donc $2n-i-1 \leq n-3$ ». C'est exactement ce
qui découle d'une annulation $H^{i}(M) = 0$ pour $i > n+1$ ; avec la borne
vraie $i > n$, on a aussi $H^{n+1}(M) = 0$, et l'application est bijective dès
$i = n+1$. Les bornes $n+2$ et $n-2$ sont surchargées et lues d'après la
remarque « $2n-i-1 \leq n-3$ » et la page 27 ; « épim » et « mono » sont des
lectures douteuses.} Ainsi $H^{q}(M)$ est connu en termes de $H^{q}(W_0)$, sauf
dans les degrés $n-1$ et $n$, où l'on a la suite exacte
\[
  H^{n-1}(M) \to H^{n-1}(W_0) \to H^{n}(M)^{\vee} \xrightarrow{\ u\ } H^{n}(M)
  \to H^{n}(W_0) \to H^{n-1}(M)^{\vee} .
\]
\footnote{La page porte « surj » sur la flèche $H^{n-1}(W_0) \to
H^{n}(M)^{\vee}$. Par exactitude, son image est le noyau de
$u : H^{n}(M)^{\vee} \to H^{n}(M)$, et elle n'est surjective que si $u$ est nul
en degré moyen, ce qui n'est pas le cas en général : pour un point double
ordinaire en $n$ pair, la forme d'intersection est la forme $(\pm 2)$ sur
$H^{n}(M) \simeq \Lambda$. On ne retient pas le mot. Les deux lignes commencent
par des termes biffés, qui reprennent la suite un cran plus bas.}

Supposons enfin que $W$ soit la sphère topologique et $n \geq 3$. Alors
$H^{i}(W_0) = 0$ pour $i \neq 0, n-1, n, 2n-1$, donc par ce qui précède
\[
  H^{i}(M) = 0 \qquad \text{pour } i \neq 0, n-1, n,
\]
et la suite des degrés du milieu se ferme en une suite exacte que la page
encadre sur deux lignes, la seconde étant la transposée de la première :
\[
  \boxed{\begin{array}{c}
    0 \to H^{n+1}(M)^{\vee} \xrightarrow{\ u\ } H^{n-1}(M) \to H^{n-1}(W_0)
      \to H^{n}(M)^{\vee} \\
    \downarrow u \\
    0 \leftarrow H^{n+1}(M) \xleftarrow{\ {}^{t}u\ } H^{n-1}(M)^{\vee}
      \leftarrow H^{n}(W_0) \leftarrow H^{n}(M)
  \end{array}}
\]
Avec $H^{n+1}(M) = 0$, les deux termes extrêmes sont nuls, et ce qui reste est
la suite à six termes
\[
  0 \to H^{n-1}(M) \to H^{n-1}(W_0) \to H^{n}(M)^{\vee} \xrightarrow{\ u\ }
  H^{n}(M) \to H^{n}(W_0) \to H^{n-1}(M)^{\vee} \to 0,
\]
symétrique sous la dualité, l'application $u$ du milieu étant
autotransposée.\footnote{L'encadré est traversé au crayon de traits en zigzag
reliant les deux lignes, avec des signes $+$ et $-$ alternés : c'est la
symétrie des deux lignes, et sans doute un calcul de caractéristique
d'Euler. On ne le reconstitue pas. Le dossier s'arrête là. Le théorème de
Milnor selon lequel $M$ a le type d'homotopie d'un bouquet de sphères de
dimension $n$ donnerait de plus $H^{n-1}(M) = 0$ pour $n \geq 2$, et
réduirait la suite à $0 \to H^{n-1}(W_0) \to H^{n}(M)^{\vee}
\xrightarrow{u} H^{n}(M) \to H^{n}(W_0) \to 0$ : la cohomologie du link est le
noyau et le conoyau de la forme d'intersection. Les pages n'atteignent pas cet
énoncé.}

\pagerange{33}{33}
\section*{En termes de fibres (page 33)}

La page 33, sur un autre papier, récrit le critère de la page 27 dans le
langage du dictionnaire de la page 16 : $V$ joue le rôle de $W$, sa partie
générique $V_{\eta}$ celui de $W - W_0$, sa partie spéciale $V_{s}$ celui de
$W_0$.\footnote{La page 16 écrit $V'$ et $V'_{s}$ ; la page 33 omet le prime.}
Il vient, pour $n \geq 1$ :
\[
  V \text{ sphère cohomologique} \iff
  \left\{\begin{aligned}
    &H^{i}(V_{\eta}) \to H^{i-1}(V_{s}) \ \text{isomorphisme pour } i \neq 0, 2n, \\
    &H^{0}(V_{\eta}) \simeq \Lambda, \\
    &H^{2n}(V_{\eta}) \to H^{2n-1}(V_{s}) \ \text{injectif, de conoyau} \simeq \Lambda .
  \end{aligned}\right.
\]
C'est le critère encadré de la page 27, rendu complet par les conditions en
degrés $0$ et $2n$. Il est accompagné de la portion de suite
\[
  0 \to H^{2n}(V_{\eta}) \to H^{2n-1}(V_{s}) \to H^{2n+1}(V) \to
  H^{2n+1}(V_{\eta}) \to H^{2n}(V_{s}) = 0,
\]
où $H^{2n}(V_{\eta}) = 0$ si $n \geq 2$ et $H^{2n+1}(V_{\eta}) = 0$ si
$n \geq 1$ — c'est l'annulation au-dessus de $n+1$, puisque $2n+1 > n+1$.

La page conclut : si l'on sait déjà que la singularité est isolée,
« $H^{i}(V_{s}) = 0$ si $i \neq 0, n-1, n$ ». Pris au pied de la lettre pour
le link, c'est faux en degré $2n-1$, où la même page a
$H^{2n-1}(V_{s}) \simeq \Lambda$. C'est en revanche exactement l'énoncé de la
page 31 pour la fibre de Milnor, et la marge, qui écrit
$H^{i}(V_{\xi}) \simeq H^{i}(V_{s})$, semble aller dans ce sens.\footnote{Doute
signalé. L'indice lu $\xi$ peut être un $\eta$ ; si c'est la fibre
géométrique générique qui est désignée, la marge et la conclusion portent sur
la fibre de Milnor. On ne corrige pas la page ; on dit ce qui est vrai :
$H^{i}(W_0) = 0$ pour $i \neq 0, n-1, n, 2n-1$, et $H^{i}(M) = 0$ pour
$i \neq 0, n-1, n$ quand $n \geq 3$.}

\pagerange{34}{34}
\section*{Un autre sujet : finitude et noyaux (page 34)}

La dernière page, sur le même papier que la précédente, change de sujet. Elle
note un « inconvénient sérieux » des groupes $\mathrm{Ext}^{n}(S, T; A)$
lorsque la catégorie $A$ satisfait l'axiome (AB5) des limites inductives
filtrantes exactes : les conditions de finitude naturelles sur les systèmes
inductifs $\mathfrak{a}$-adiques, ou essentiellement $\mathfrak{a}$-adiques —
pour $A$ localement noethérienne — ne semblent pas stables par passage au
noyau. Dualement, sous (AB5$^{*}$), les conditions de finitude correspondantes
ne semblent pas stables par conoyau. La phrase finit sur deux points
d'exclamation.\footnote{Le premier paragraphe est le plus difficile de la page :
le mot qui précède « Ext » est illisible, « Ext » et « noyaux » sont des
lectures douteuses, et l'on ne précise pas davantage ce que sont $S$ et $T$,
que la page ne définit pas. Ce que l'on appellerait aujourd'hui ces systèmes
inductifs est le langage des ind-objets et des modules de torsion
$\mathfrak{a}$-primaire.}

Suit un \emph{Exemple}, entièrement barré : $A$ un anneau noethérien, $f$, $g$
deux éléments formant une suite $M$-régulière, et
\[
  P = M_{f}/M = \varinjlim_{n} M/f^{n}M,
\]
le système inductif ayant pour flèches de transition $M/f^{n}M \to M/f^{m}M$
la multiplication par $f^{m-n}$, donc injectives. L'exemple s'arrête sur
« $Q =$ ».\footnote{L'objet $M_{f}/M$ est le premier module de cohomologie
locale $H^{1}_{(f)}(M)$ ; la régularité de la suite $(f, g)$ suggère que le
contre-exemple annoncé devait se construire à partir de lui et de la
multiplication par $g$, mais la page ne le dit pas, et on ne le complète pas.
La lettre $P$ est ici sans rapport avec la variété à bord des pages
précédentes.}

\end{document}
